%%%%%%%%%%%%%%%%%%%%%%%%%%%%%%%%%%%%%%%%%
% Beamer Presentation
% LaTeX Template
% Version 1.0 (10/11/12)
%
% This template has been downloaded from:
% http://www.LaTeXTemplates.com
%
% License:
% CC BY-NC-SA 3.0 (http://creativecommons.org/licenses/by-nc-sa/3.0/)
%
%%%%%%%%%%%%%%%%%%%%%%%%%%%%%%%%%%%%%%%%%

%----------------------------------------------------------------------------------------
%	PACKAGES AND THEMES
%----------------------------------------------------------------------------------------

\documentclass{beamer}

\mode<presentation> {

% The Beamer class comes with a number of default slide themes
% which change the colors and layouts of slides. Below this is a list
% of all the themes, uncomment each in turn to see what they look like.

%\usetheme{default}
%\usetheme{AnnArbor}
%\usetheme{Antibes}
%\usetheme{Bergen}
%\usetheme{Berkeley}
%\usetheme{Berlin}
%\usetheme{Boadilla}
%\usetheme{CambridgeUS}
%\usetheme{Copenhagen}
%\usetheme{Darmstadt}
%\usetheme{Dresden}
%\usetheme{Frankfurt}
%\usetheme{Goettingen}
%\usetheme{Hannover}
%\usetheme{Ilmenau}
%\usetheme{JuanLesPins}
%\usetheme{Luebeck}
%\usetheme{Madrid}
%\usetheme{Malmoe}
%\usetheme{Marburg}
%\usetheme{Montpellier}
%\usetheme{PaloAlto}
\usetheme{Pittsburgh}
%\usetheme{Rochester}
%\usetheme{Singapore}
%\usetheme{Szeged}
%\usetheme{Warsaw}

% As well as themes, the Beamer class has a number of color themes
% for any slide theme. Uncomment each of these in turn to see how it
% changes the colors of your current slide theme.

%\usecolortheme{albatross}
\usecolortheme{beaver}
%\usecolortheme{beetle}
%\usecolortheme{crane}
%\usecolortheme{dolphin}
%\usecolortheme{dove}
%\usecolortheme{fly}
%\usecolortheme{lily}
%\usecolortheme{orchid}
%\usecolortheme{rose}
%\usecolortheme{seagull}
%\usecolortheme{seahorse}
%\usecolortheme{whale}
%\usecolortheme{wolverine}

%\setbeamertemplate{footline} % To remove the footer line in all slides uncomment this line
%\setbeamertemplate{footline}[page number] % To replace the footer line in all slides with a simple slide count uncomment this line

%\setbeamertemplate{navigation symbols}{} % To remove the navigation symbols from the bottom of all slides uncomment this line
}

\usepackage{graphicx} % Allows including images
\usepackage{booktabs} % Allows the use of \toprule, \midrule and \bottomrule in tables

\usepackage{fontspec}
\setsansfont{Calibri}[
    Path=Font/,
    Scale=0.9,
    Extension = .ttf,
    UprightFont=*-Regular,
    BoldFont=*-Bold,
    ItalicFont=*-Italic,
    BoldItalicFont=*-BoldItalic
    ]
\usepackage{unicode-math}
\setmathfont{Libertinus Math}% 使用默认的数学字体或你偏好的数学字体 Latin Modern Math/Libertinus Math

\usepackage{ragged2e}
\justifying\let\raggedright\justifying %两端对齐

 \usepackage{tcolorbox} 
%----------------------------------------------------------------------------------------
%	TITLE PAGE
%----------------------------------------------------------------------------------------

\title[Asset Pricing]{\textbf{Asset Pricing}} % The short title appears at the bottom of every slide, the full title is only on the title page

\subtitle{Why Focus on SDF Innovation?}

\author[Cheng Hang]{Cheng Hang} % Your name
\institute[DUFE] % Your institution as it will appear on the bottom of every slide, may be shorthand to save space
{School of Fintech \\
Dongbei University of Finance \& Economics \\ % Your institution for the title page
\medskip
\textbf{chenghangch@qq.com}% Your email address
}
\date{\today} % Date, can be changed to a custom date

	\AtBeginSection[]
{
	\begin{frame}{Content}
		\tableofcontents[sectionstyle=show/shaded,subsectionstyle=show/shaded/hide] %这几个参数我也不知道该如何准确地解释，反正它们最终的效果是突出显示当前章节，而其它章节都进行了淡化处理
		\addtocounter{framenumber}{-1}  %目录页不计算页码
	\end{frame}
}

%----------------------------------------------------------------------------------------
%	Begin
%----------------------------------------------------------------------------------------

\begin{document}

\begin{frame}
\titlepage % Print the title page as the first slide
\end{frame}


\begin{frame}[allowframebreaks]
    \frametitle{Why Focus on SDF Innovation?}
    
    \textbf{Question:} Why do we decompose the SDF into expected and unexpected components?
    
    \textbf{Recall the fundamental asset pricing equation:}
    \begin{equation*}
        1 = E_t[M_{t+1} R_{i,t+1}]
    \end{equation*}
    
    Expanding this:
    \begin{equation*}
        1 = E_t[M_{t+1}] \cdot E_t[R_{i,t+1}] + \text{Cov}_t(M_{t+1},R_{i,t+1})
    \end{equation*}
    
    \framebreak
    
    \textbf{Key insight:} The expected component $E_t[M_{t+1}]$ is common to all assets.
    
    Rearranging for expected return:
    \begin{equation*}
        E_t[R_{i,t+1}] = \frac{1}{E_t[M_{t+1}]}  - \frac{\text{Cov}_t(M_{t+1}, R_{i,t+1})}{E_t[M_{t+1}]}
    \end{equation*}
    
    \textbf{The risk-free rate:}
    \begin{equation*}
      R_{f,t} = \frac{1}{E_t[M_{t+1}]}
    \end{equation*}
    
    \textbf{Therefore, the risk premium is:}
    \begin{tcolorbox}[colback=orange!5!white, colframe=orange!75!black]
    \begin{equation*}
        E_t[R_{i,t+1}] - R_{f,t} = -\frac{\text{Cov}_t(M_{t+1}, R_{i,t+1})}{E_t[M_{t+1}]}
    \end{equation*}
    \end{tcolorbox}
    
    \framebreak
    
    \textbf{Observation:} Risk premium depends ONLY on the covariance of $M_{t+1}$ with returns.
    
    Since $E_t[M_{t+1}]$ is known at time $t$, we can write:
    \begin{equation*}
        \text{Cov}_t(M_{t+1}, R_{i,t+1}) = \text{Cov}_t(M_{t+1} - E_t[M_{t+1}], R_{i,t+1}) = \text{Cov}_t(\tilde{M}_{t+1}, R_{i,t+1})
    \end{equation*}
    
    where $\tilde{M}_{t+1} \equiv M_{t+1} - E_t[M_{t+1}]$ is the SDF innovation.
    
    \textbf{Conclusion:} For asset pricing (risk premia), only the unexpected component of the SDF matters!
    
    \framebreak
    
    \textbf{Economic interpretation in three points:}
    
    \textbf{1. Risk premium is about unexpected shocks}
    \begin{itemize}
        \item Assets are risky because their payoffs covary with unexpected changes in marginal utility
        \item Expected changes in marginal utility are already reflected in the risk-free rate
        \item Only surprises create risk that requires compensation
    \end{itemize}
    
    \textbf{2. SDF innovation identifies priced risk factors}
    \begin{itemize}
        \item Each term in $\tilde{m}_{t+1}$ represents a source of systematic risk
        \item Assets with high covariance with these shocks command risk premia
        \item This is the foundation for multifactor models
    \end{itemize}
    
    \framebreak
    
    \textbf{3. SDF innovation reveals what investors fear}
    \begin{itemize}
        \item The terms that appear in $\tilde{m}_{t+1}$ tell us what kinds of shocks affect marginal utility
        \item The coefficients tell us how much investors care about each shock
        \item This guides us in building consumption-based asset pricing models
    \end{itemize}
    
    \framebreak
    
    \textbf{Log-linear approximation:}
    
    Under joint lognormality, we can work with log returns and log SDF:
    \begin{equation*}
        E_t[r_{i,t+1}] - r_{f,t} \approx -\text{Cov}_t(\tilde{m}_{t+1}, r_{i,t+1}) - \frac{1}{2}\text{Var}_t(r_{i,t+1})
    \end{equation*}
    
    where lowercase letters denote logs.
    
    \textbf{This is why we focus on the log-linearized SDF innovation:}
    \begin{equation*}
        \tilde{m}_{t+1} = m_{t+1} - E_t[m_{t+1}]
    \end{equation*}
    
    It directly determines risk premia through its covariance with asset returns.
\end{frame}

\end{document}